% CESC 410L lab report.
%
% NOTE TO WHOEVER EDITS THIS: keep this file free of references to any
% tooling, scripts or repository paths. It is a submitted document. Build and
% packaging instructions belong in reference_docs/report_guide.md, not here.
% Comment-only lines are stripped at package time as a backstop, but do not
% rely on that.

\documentclass[11pt]{article}

\usepackage[margin=1in]{geometry}
\usepackage{amsmath, amssymb}     % equations
\usepackage{graphicx}             % \includegraphics
\usepackage{booktabs}             % clean tables
\usepackage{listings}             % code snippets
\usepackage{xcolor}
\usepackage[colorlinks=true, allcolors=blue]{hyperref}
\usepackage{caption}

% Code style: Courier New equivalent at 10pt, per the CEC 320 template's
% "code in Courier New 12, single spacing". Light background, no dark blocks --
% reports may be printed.
\lstset{
  basicstyle=\ttfamily\small,
  breaklines=true,
  frame=single,
  backgroundcolor=\color{gray!8},
  showstringspaces=false,
  captionpos=b,
  language=Python,
}

\title{CESC 410L --- Lab NN: \textit{Title}}
\author{Name(s) \\ Section NN}
\date{Lab start: MM/DD/YYYY \\ Report date: MM/DD/YYYY}

\begin{document}
\maketitle

%======================================================================
\section{Introduction}
%======================================================================

What the lab asks for, in your own words, and why. Include any preparation
or background reading. Keep it short.

%======================================================================
\section{Narrative}
%======================================================================

How the procedure differed for you: modifications, problems, bugs, and how
they were resolved. Research done to fix something belongs here.

This is where deviations from the handout go. It is a graded section, and a
real debugging story is exactly what it asks for --- take it from the
\textbf{Deviations} table in \texttt{labNN/README.md}.

%======================================================================
\section{Artifacts}
%======================================================================

Only code you added, modified, or commented on. Name the source file above
each snippet.

\begin{lstlisting}[caption={\texttt{sinusoids.py} --- example}]
def plot_one_sinusoid_td_fd(Am, Fr, Ph, ...):
    # Take the DFT of the sampled signal, shifting zero frequency
    # to the centre so the spectrum plots symmetrically about 0 Hz.
    X_s = np.fft.fft(x_s)
    X_F = np.abs(np.fft.fftshift(X_s))
\end{lstlisting}

\begin{figure}[h]
  \centering
  \includegraphics[width=0.8\textwidth]{figs/lab0_sinusoids_fig1.png}
  \caption{Single sinusoid, $A=1$, $f=4$\,Hz, $\phi=\pi/4$. Sparse samples
           (dots) over the dense waveform (line), with the magnitude
           spectrum below.}
  \label{fig:one-sinusoid}
\end{figure}

Refer to figures by number, as in Figure~\ref{fig:one-sinusoid}. Avoid dark
backgrounds in screenshots --- reports may be printed.

%======================================================================
\section{Discussion and Results}
%======================================================================

What you learned and how it applies later. Then, explicitly:

\subsection*{Handout questions}

Handouts bury questions between code blocks and they are graded. Answer each
one here.

\begin{enumerate}
  \item \textbf{Q.} \textit{question} \\
        \textbf{A.} answer
\end{enumerate}

\subsection*{Results}

Summarise. Where a result rests on an equation, state it --- this is why the
report is in \LaTeX{}. For a sampled sinusoid
%
\begin{equation}
  x[n] = A\cos\!\left(2\pi \frac{f}{f_s} n + \phi\right),
  \qquad n = 0,\dots,N-1,
  \label{eq:sampled-sinusoid}
\end{equation}
%
the discrete Fourier transform is
%
\begin{equation}
  X[k] = \sum_{n=0}^{N-1} x[n]\, e^{-j 2\pi k n / N},
  \qquad k = 0,\dots,N-1,
  \label{eq:dft}
\end{equation}
%
so a real sinusoid whose frequency falls on a bin centre produces two lines,
at $\pm f$, each of magnitude $AN/2$.

\begin{table}[h]
  \centering
  \caption{Signal parameters.}
  \begin{tabular}{lrrr}
    \toprule
    Component & $A$ & $f$ (Hz) & $\phi$ (rad) \\
    \midrule
    1 & 1.0 & 4  & $\pi/4$ \\
    2 & 0.5 & 8  & $\pi/3$ \\
    3 & 0.2 & 12 & $\pi/9$ \\
    \bottomrule
  \end{tabular}
\end{table}

\end{document}
